% SEGMENT OLP-0579-S01
% Part: set-theory
% Chapter: ord-arithmetic
% Section: exponentiation

\documentclass[../../../include/open-logic-section]{subfiles}

\begin{document}

\olfileid{sth}{ord-arithmetic}{expo}
\olsection{வரிசையெண் அடுக்கேற்றம்}

இப்போது வரிசையெண் அடுக்கேற்றத்தை நோக்குவோம். வருத்தமாக,
இதற்கான \emph{எளிதில் புரியும்} கட்டமைப்பு வழி வரையறை இல்லை.

வெளிப்படையான கட்டமைப்பு வழி வரையறைகள் \emph{உள்ளன} என்பது உண்மை.
அவற்றில் ஒன்று இதோ. $\text{finfun}(\alpha,\beta)$ என்பது
$f \colon \alpha \to \beta$ என்ற எல்லாச் சார்புகளிலும்,
$\Setabs{\gamma \in \alpha}{f(\gamma) \neq 0}$ ஏதோ ஓர் இயல்
எண்ணுடன் எண்ணொத்ததாக உள்ள சார்புகளின் கணம் என்க.
$\text{finfun}(\alpha,\beta)$ மீது, $f \neq g$ என்றும்
$f(\gamma_0) < g(\gamma_0)$ என்றும் இருந்தால், அப்போது மட்டுமே
$f \sqsubset g$ என நன்கு வரிசைப்படுத்தலை வரையறுக்கிறோம்; இங்கு
$\gamma_0 = \text{max}\Setabs{\gamma \in \alpha}{f(\gamma) \neq
g(\gamma)}$. அப்போது $\ordexpo{\alpha}{\beta}$ வைப்
$\ordtype{\text{finfun}(\beta, \alpha), \sqsubset}$ என
வரையறுக்கலாம்.
% SOURCE CORRECTION TA-STH-028: exponent beta indexes finite-support functions into base alpha.
பாட்டர் இந்த வெளிப்படையான வரையறையைப் பயன்படுத்தியவுடன் இவ்வாறு
விளக்குகிறார்:
\begin{quote}
  இந்த வரிசையைத் தேர்ந்தெடுப்பது, பொருத்தமான மீள்வரையறை
  நிபந்தனையைப் பூர்த்திசெய்யும் வரிசையெண் அடுக்கேற்ற வரையறையைப்
  பெற வேண்டும் என்ற நமது விருப்பத்தால் மட்டுமே தீர்மானிக்கப்படுகிறது
  \ldots; வரிசைப்படுத்தப்பட்ட கூட்டுத்தொகையையும் பெருக்கற்பலனையும்
  விட இதை மனக்கண்ணில் காண்பது மிகவும் கடினம்.
  \citep[p.~199]{Potter2004}
\end{quote}
அப்படித்தான். கூட்டலை ``$\alpha$ வின் பிரதிக்குப் பிறகு
$\beta$ வின் பிரதி'' என்றும், பெருக்கலை ``$\alpha$ வின்
பிரதிகளால் ஆன $\beta$-தொடர்'' என்றும் விளக்கினோம். ஆனால்
$\text{finfun}(\alpha, \gamma)$ பற்றி அத்தகைய சுருக்கமான
விளக்கம் நம்மிடம் இல்லை. எனவே வரிசையெண் அடுக்கேற்றத்தை
முடிவிலிக்கடந்த மீள்வரையறை \emph{மூலமே} வரையறுப்போம்; அதாவது:

% SEGMENT OLP-0579-S02
\begin{defn}\ollabel{ordexporecursion}
\begin{align*}
  \ordexpo{\alpha}{0} &= 1\\
  \ordexpo{\alpha}{\beta\ordplus 1} &=\ordexpo{\alpha}{\beta} \ordtimes \alpha\\
  \ordexpo{\alpha}{\beta} &= \bigcup_{0 < \delta < \beta}\ordexpo{\alpha}{\delta}& & \text{$\beta$ ஓர் எல்லை வரிசையெண் எனில்}
\end{align*}
% SOURCE CORRECTION TA-STH-029: omit the zero exponent at positive limits to keep 0^beta=0.
\end{defn}

\emph{கணக் கோட்பாட்டாளர்களாகவே} நாம் செயல்பட்டிருந்தால், வரிசையெண்
அடுக்கேற்றத்தின் சில பண்புகளை ஆராய விரும்பியிருக்கலாம். ஆனால்
வரிசையெண் அடுக்கேற்றம் பரிமாற்றுப் பண்புடையதல்ல என்ற
ஆச்சரியமில்லாத உண்மையைத் தவிர வேறொன்றையும் இங்கு பெரிதாகச்
சேர்க்க வேண்டியதில்லை. ஆக,
$\ordexpo{2}{\omega} =
\bigcup_{\delta < \omega}\ordexpo{2}{\delta} = \omega$;
ஆனால் $\ordexpo{\omega}{2} = \omega \ordtimes \omega$.
இயல் எண்களிலேயே அடுக்கேற்றம் பரிமாற்றுப் பண்புடையதல்ல என்பதால்
அப்படியிருக்க வேண்டும் என்று \emph{எதிர்பார்க்கவும் கூடாது}:
$\ordexpo{2}{3} = 8 < 9 = \ordexpo{3}{2}$.

\begin{prob}
முடிவிலிக்கடந்த தொகுத்தறிதலைப் பயன்படுத்தி,
$\ordexpo{\alpha}{\beta} =
\ordtype{\text{finfun}(\beta, \alpha), \sqsubset}$ என
வரையறுத்தால், \olref[sth][ord-arithmetic][expo]{ordexporecursion}
இன் மீள்வரையறைச் சமன்பாடுகள் கிடைக்கின்றன என்பதை நிறுவுக.
% SOURCE CORRECTION TA-STH-028: the exercise uses the corrected argument order.
\end{prob}

\end{document}
